\begin{trapbox}
\begin{description}[style=nextline, leftmargin=2.8em, labelsep=0.5em, itemsep=0.35em]
    \item[\textbf{\textcolor{CTrap}{易错点一（混淆旋转轴）}}] 乱用公式 $V=\frac{4}{3}\pi a^2 b$ 和 $V=\frac{4}{3}\pi a b^2$，不考虑旋转轴。
    \item[\textbf{\textcolor{CTrap}{正确理解（轴与平方对应）：}}] 绕 $y$ 轴时 $a$（$x$ 轴半轴）平方，绕 $x$ 轴时 $b$（$y$ 轴半轴）平方。
    \item[\textbf{\textcolor{CTrap}{错因分析（记忆方法）：}}] 死记硬背未理解来源。通过积分推导可发现平方永远出现在垂直于旋转轴的半轴上。
\end{description}

\medskip
\begin{description}[style=nextline, leftmargin=2.8em, labelsep=0.5em, itemsep=0.35em]
    \item[\textbf{\textcolor{CTrap}{易错点二（半轴长短混淆）}}] 忽视 $a$、$b$ 与坐标轴的对应关系，误认为 $a$ 是长半轴、$b$ 是短半轴并代入公式。
    \item[\textbf{\textcolor{CTrap}{正确理解（坐标定义明确）：}}] $a$ 始终是 $x$ 轴方向半轴长，$b$ 是 $y$ 轴方向半轴长，无论长短。
    \item[\textbf{\textcolor{CTrap}{错因分析（习惯固化）：}}] 椭圆通常有长、短半轴概念，但体积公式直接使用方程参数，与长短无关。
\end{description}

\medskip
\begin{description}[style=nextline, leftmargin=2.8em, labelsep=0.5em, itemsep=0.35em]
    \item[\textbf{\textcolor{CTrap}{易错点三（系数记错）}}] 将 $V_y$ 记成 $\frac{4}{3}\pi a b^2$ 或 $\frac{4}{3}\pi a b$ 等。
    \item[\textbf{\textcolor{CTrap}{正确理解（公式系数精准）：}}] $V_y = \frac{4}{3}\pi a^2 b$，系数 $\frac{4}{3}$，$\pi$，$a^2 b$ 缺一不可。
    \item[\textbf{\textcolor{CTrap}{错因分析（类比不当）：}}] 常与球体积 $\frac{4}{3}\pi R^3$ 混淆，需通过积分或记忆结构区分。
\end{description}
\end{trapbox}
